% !TEX program = xelatex
% Standalone, original figures inline. No parent aux or external TeX required.
\documentclass[11pt,UTF8,fontset=fandol]{ctexart}
\usepackage[a4paper,margin=24mm,headheight=16pt]{geometry}
\usepackage{amsmath,amssymb,amsthm,mathtools,bm,booktabs,tabularx,longtable,array}
\usepackage{tikz,xcolor,enumitem,fancyhdr,graphicx}
\usetikzlibrary{arrows.meta,positioning,calc,fit,matrix}
\usepackage[most]{tcolorbox}
\usepackage[section]{placeins}
\usepackage[unicode,bookmarksnumbered,colorlinks=true]{hyperref}
\definecolor{ink}{RGB}{24,53,77}
\definecolor{teal}{RGB}{26,120,113}
\definecolor{amber}{RGB}{173,102,35}
\definecolor{purple}{RGB}{107,75,135}
\hypersetup{linkcolor=ink,urlcolor=teal,citecolor=teal,
 pdftitle={TabU-TAR：面向实现的模型设计},pdfauthor={TabU Project}}
\IfFontExistsTF{lmroman10-regular.otf}{\setmainfont{lmroman10-regular.otf}[
 BoldFont=lmroman10-bold.otf,ItalicFont=lmroman10-italic.otf,
 BoldItalicFont=lmroman10-bolditalic.otf]}{}
\ctexset{section={format=\Large\bfseries\color{ink}},subsection={format=\large\bfseries}}
\setlength{\parskip}{0.3em}
\setlength{\emergencystretch}{3em}
\setlist{leftmargin=2em,itemsep=0.2em}
\renewcommand{\arraystretch}{1.2}
\newcolumntype{Y}{>{\raggedright\arraybackslash}X}
\pagestyle{fancy}\fancyhf{}
\fancyhead[L]{\small TabU-TAR / Model Design}\fancyhead[R]{\small Proposed realization}
\fancyfoot[C]{\thepage}
\setcounter{tocdepth}{1}\numberwithin{equation}{section}\allowdisplaybreaks
\newcommand{\R}{\mathbb R}\newcommand{\V}{\mathcal V}\newcommand{\Q}{\mathcal Q}
\newcommand{\Nul}{\mathcal N}\newcommand{\ind}{\mathbf 1}
\newcommand{\sm}{\operatorname{softmax}}\newcommand{\OMAB}{\operatorname{OMAB}}
\newcommand{\Norm}{\operatorname{RMS}}\newcommand{\FFN}{\operatorname{FFN}}
\newtheorem{proposition}{命题}[section]\renewcommand{\proofname}{证明}
\newtcolorbox{decision}{colback=teal!4,colframe=teal!65!black,boxrule=0.5pt,arc=1mm}
\tikzset{box/.style={draw=ink,fill=ink!3,rounded corners=2pt,align=center,
 inner sep=5pt,font=\small},fact/.style={box,draw=teal,fill=teal!7},
 choice/.style={box,draw=amber,fill=amber!7},truth/.style={box,draw=purple,dashed,fill=purple!5},
 arr/.style={-{Latex[length=2mm]},semithick,draw=ink}}
\begin{document}
\hypersetup{pageanchor=false}
\begin{titlepage}
\vspace*{17mm}
{\color{teal}\large TabU / Typed Additive Readout}\par\vspace{8mm}
{\Huge\bfseries\color{ink} TabU-TAR}\par\vspace{5mm}
{\LARGE 面向实现的模型设计}\par\vspace{4mm}
{\large 以 TabPFN-3 为规模参照的首份完整配置}\par\vspace{12mm}
\begin{decision}
\textbf{54,071,520 个持久参数}\quad $d=384$，$L=15$，$K=32$，$S=128$。\\
单轴固定 inducing slots；逐 Cell 表示；类型化加性响应；同列 LL/NW 预测。\\
保留 without inducing baseline，并分别给出同深度与近似同参数量配置。
\end{decision}
\vspace{8mm}
\textbf{这份稿件的用途}\par
定义一个可以被实现、检查、训练和比较的具体模型。阅读本文件即可获得输入合同、
全部前向公式、张量形状、参数共享、初始化、训练起始配方、推理边界及实现验收条件。
它也作为未来论文的首份骨架：方法部分完整，实验问题已经提出，结果尚待产生。
\par\vspace{4mm}
\textbf{状态}\par
设计提案，2026 年 9 月 6 日。参数量由张量清单计算，尚未构建或训练该模型。
本文没有性能结果，不承诺达到 TabPFN-3 的精度、容量或速度。
“首份配置”不构成新 runtime 版本号，也不绑定任何旧 checkpoint。
\vfill
\small
共同数学由 \texttt{MATHEMATICAL\_DESIGN.tex} 维护；本稿选择其中一组 realization。
\texttt{design/defaults.json} 提供机器可读数值；本稿的数学定义是实现语义依据。
所有图均为本项目绘制的机制示意，不是实验结果图。
\end{titlepage}
\hypersetup{pageanchor=true}

\begin{abstract}
TabU-TAR 将表格预测表达为：从可见 Cell 构造证据，在表格轴向传播中形成 Cell、Unit、
Feature 三类表示，再通过共享的类型化线性读出得到加性响应，并以同列可见样本形成数值或离散预测。
本文把这一结构具体化为约 54M 参数的模型，参考 TabPFN-3 的公开规模与若干工程默认，
采用固定 128 槽的列内 inducing 通信和直接行内交互。与压缩为单一行表示的路线相比，
本设计保留逐 Cell 查询接口，同时承担更大的二维激活成本。
本文完整定义前向、局部线性求解、训练梯度和失败状态；指出默认 terminal 对 Feature 公共平移不敏感，
并把这一限制作为明确的模型取舍。所列训练配方用于开始实现与验证，尚不是经过实验选定的预训练方案。
\end{abstract}
\tableofcontents
\clearpage

\section{设计目标与参照边界}
\label{sec:position}

目标是得到一个\textbf{单一、具体、可复现的通用表格模型设计}，能够处理 numeric、nominal、ordinal
列的 masked-cell recovery，并通过把目标列的待预测值标成 Query 表达监督预测。
本文不给 RecSys、Causal、Graph 引入场景专属算子；这些设计仍可在共同接口上独立生长。

TabPFN-3 的报告 v2 给出分类约 53M、回归约 58M 参数；其附录 C 明确列出发布模型配置\cite{pfn3}。
TAR 以这一区间作为\emph{持久参数预算}，没有复制其整套结构，也没有继承其验证过的数据规模。
表~\ref{tab:reference} 是本稿使用的对照，具体值以报告附录而非可变的 GitHub 默认配置为准。

\begin{table}[htbp]\centering\small
\caption{规模与少量结构默认的对照。右列是本文的设计选择。}\label{tab:reference}
\begin{tabularx}{\linewidth}{@{}lYY@{}}\toprule
维度 & TabPFN-3 报告 v2 & 本文 TAR \\\midrule
持久参数 & 分类约 53M；回归约 58M & 54.071520M，同一 typed 模型 \\
载体宽度 & 列编码 128；行 ICL 512 & 全部 Cell/bank/slot 均为 384 \\
主要深度 & 列 3；聚合 3；ICL 24 & 15 个列后行 block \\
Attention heads & 列、行与 ICL 均 8 & 8，每头 48 维 \\
FFN expansion & 2 & 2，即 768 \\
Inducing 数 & 128 & 128，仅列内样本轴 \\
行汇聚 & 4 个 CLS 拼接 & 保留 Cell 与 $K=32$ 语义 banks \\
特征分组 & 循环三元组 & 单 Cell，不做三元组分组 \\
Normalization & RMSNorm & FFN 前逐 token RMSNorm \\
输出 & 分类 retrieval；回归分桶分布 & 离散平滑 NW；数值 LL 点预测 \\\bottomrule
\end{tabularx}
\end{table}

参考实现的固定 commit 和报告原稿位置见 \texttt{references/tabpfn-3/}。
代码用于核对 inducing 的收集/回读与小样本路径\cite{pfncode}，不被当作 TAR 的运行依赖。
优化器、TAR 宽度/深度、$K$、初始化与训练表分布均为\textbf{本文新选的起始默认}，不是复述对方训练配方。

\begin{decision}
TAR 当前默认的理由：用 $d=384$ 降低逐 Cell 激活宽度，以 15 个 block 达到约 54M 参数；
用固定 $S=128$ 控制列内通信；用 $K=32$ 限制语义 bank 和 LL 求解成本。
这些是预算与接口上的选择，其效果必须通过训练验证。
\end{decision}

\subsection{面向未来论文的中心问题}
方法贡献应围绕可检查的问题组织：逐 Cell 的类型化接口是否支持跨表共享学习；
固定槽位是否在压缩成本时保留有用证据；局部 typed estimation 是否能利用学习到的响应几何。
这三项目前是研究问题。数学上的零元闭合、Query 插入不扰动、LL 唯一性，不能替代预测质量证据。
后文将完整定义方法；第~\ref{sec:experiments} 节仅给待完成的实验结构。

\section{一张表如何成为一次任务}
\label{sec:input}

输入是 $N$ 个 Unit（行）和 $M$ 个 Feature（列），$N,M$ 均为正整数。
地址须唯一且位于声明范围内，schema 按列完整给定。所有下标均是 episode-local 地址，
不是持久 learned ID。列 schema 声明：numeric 的答案域为 $\R$；nominal 为非空有限集合
$\mathcal D_a$；ordinal 另带声明的全序。文本、日期需上游显式转成这些类型；本稿不定义文本编码器。

\subsection{输入、可见事实和真值的分离}
定义自然观测地址 $\mathcal O$、待回答地址 $\Q$、可见事实 $\V$，要求
$\V\cap\Q=\varnothing$。训练时 $\Q\subseteq\mathcal O$ 且 $\V=\mathcal O\setminus\Q$；
推理 Query 可以是原来未知的地址。其他原始 Cell 是 Null。
Forward 只接收
\begin{equation}
 E_{\rm in}=(N,M,\text{schema},\V,\{x_{ra}:(r,a)\in\V\},\Q,\xi_{\rm prep}).
 \label{eq:input}
\end{equation}
隐藏值单独存于 $Y^\star_{\Q}$，仅 scorer 读取。输入 tensor 的隐藏位置不能继续存放真值后仅靠
attention mask 遮挡；应在编译前移除，避免统计量、类别字典或预处理读取它。

\begin{figure}[htbp]\centering
\begin{tikzpicture}[node distance=7mm and 8mm]
\node[fact,text width=39mm] (e) {可见值、schema、地址\\$E_{\rm in}$};
\node[box,text width=39mm,right=of e] (c) {类型预处理与 compiler\\$H^{(0)}$};
\node[choice,text width=39mm,right=of c] (b) {15 个轴向 block\\$H^{(L)}$};
\node[box,text width=39mm,below=13mm of b] (r) {Cell / Unit / Feature\\加性响应 $z$};
\node[fact,text width=39mm,left=of r] (t) {同列可见支持\\typed LL / NW};
\node[box,text width=39mm,left=of t] (p) {点预测或 PMF\\以及证据状态};
\node[truth,text width=39mm,below=10mm of p] (y) {隐藏 Query 真值\\$Y^\star_{\Q}$};
\node[truth,text width=39mm,right=of y] (loss) {训练时唯一汇合点\\scorer / loss};
\draw[arr] (e)--(c);\draw[arr] (c)--(b);\draw[arr] (b)--(r);
\draw[arr] (r)--(t);\draw[arr] (t)--(p);\draw[arr] (p.south)--++(0,-3mm)-|(loss.north);
\draw[arr] (y)--(loss);
\end{tikzpicture}
\caption{端到端信息流。绿色是事实及其聚合，橙色是有参数的槽位通信，虚线框是训练边界。
支持值从 $E_{\rm in}$ 读取；真值不参与任何上游计算。}
\label{fig:flow}
\end{figure}

\subsection{证据合同与 API 状态}
每个目标列的支持地址为 $\mathcal C_a=\{s:(s,a)\in\V\}$，$n_a=|\mathcal C_a|$。
默认 $n_a\ge2$ 才输出预测。$n_a=0,1$ 对该目标返回 \texttt{insufficient-support}，不自动改用先验或单样本答案；
同一输入中的其他合法目标可继续返回预测。
非法类型、非有限 numeric visible value、域外离散值、重叠 masks 返回 \texttt{invalid-input}。
自然 NaN/Inf 在入口被标为未观测；不得将其视为值零。
实现内部的非有限激活或求解失败返回 \texttt{numerical-failure}，不输出伪造的有限预测。

分类声明域不得从隐藏 labels 推断。监督分类默认使用训练标签出现的类别作为声明域；
测试中出现域外类别时，评估明确记录 out-of-domain，不能偷偷扩充输入 schema。
若任务方提前提供更大的合法域，可以包括未见类别；这些类别只有平滑质量，不凭空产生事实。

\subsection{监督预测的适配}
把 $(X,y)$ 拼成一张表，目标标签列在训练行可见、测试行是 Query。
\texttt{predict} 的初始默认为\textbf{逐测试行的独立 episode}：全部训练行加一条测试行，
测试行已知 covariates 也属于 $\V$。不同测试行分别调用，返回顺序按外部 ID 恢复。
因此单个预测不依赖其他测试行是否被同批提交；这不意味着测试行自身 covariates 不影响支持表示。

\texttt{predict\_cells} 则可在同一明确的 $E_{\rm in}$ 中回答多个 Cell。
多测试行联合放入 $\V$ 是另一种 transductive 配置，会改变证据与预测；不能把它当作等价的 batching 优化。
这一点与后文缓存条件直接相关。

\section{编码：共享表示，不保存行列身份}
\label{sec:encoding}

\subsection{Numeric 与 ordinal 的连续通道}
Numeric 仅用同列可见值计算总体均值 $\widehat\mu_a$ 与总体方差 $\widehat\sigma_a^2$：
\begin{equation}
 t_{ra}=\frac{x_{ra}-\widehat\mu_a}{\widehat\sigma_a+10^{-6}}.
 \label{eq:numeric-prep}
\end{equation}
常数列得到 $t=0$；无可见值列没有 Value 需要编码。默认不裁剪、不加入 SVD、不使用 quantile ensemble。
Ordinal 若域大小为 $D_a>1$，取 $t=(\operatorname{rank}(x)-1)/(D_a-1)$；$D_a=1$ 时取零。
Numeric/ordinal 共用一个 continuous branch，没有额外 learned type embedding。

取 $J=d/4=96$，可学习频率 $\omega\in\R^J$、无 bias 投影 $A\in\R^{d\times2J}$：
\begin{equation}
 \phi_\omega(t)=[\sin(\omega_1t),\ldots,\sin(\omega_Jt),
 \cos(\omega_1t),\ldots,\cos(\omega_Jt)]^\top,
 \qquad h^{\rm val}_{ra}=A\phi_\omega(t_{ra}).
 \label{eq:fourier}
\end{equation}
$\omega,A$ 跨列、类型与 episode 共享。频率不按列重新拟合。完整初始化见第~\ref{sec:init} 节。
Fourier lift 恒非零不等于无约束训练后的投影恒非零；实现需检查非 Null 初始 carrier 的有限非零性。

\subsection{Nominal 通道与可复现 codebook}
对每列的每个\emph{可见类别}抽取 $g_{a,b}\sim\mathcal N(0,I_d)$，令
\begin{equation}
 e_{a,b}=g_{a,b}/\|g_{a,b}\|_2,\qquad h^{\rm val}_{ra}=e_{a,x_{ra}}.
 \label{eq:nominal}
\end{equation}
该通道没有 learned category embedding 或额外投影；同列同类别共享 code，不同列独立。
未见类别只占输出域坐标，不进入 codebook。每个新训练 episode 重抽 codebook；重放、分块和比较时固定。
保存实际 codebook 可消除不同 RNG 后端的差异。
实现可按 schema 类别列表的固定编号枚举可见类别，再按稳定列 ID 分配 RNG 子流；
该列表顺序也用于输出 PMF 坐标和 argmax 并列处理。导入后的 ID 不随 tensor 重排重新编号。

\clearpage
分配随机数必须绑定稳定的 episode、列与类别外部标识，不能绑定当前 tensor 遍历序号。
在重排行列或重标类别的等变性检查中，codebook 与对应身份一起移动；重新采样只保证分布意义的对称性。
Gaussian draw 的零/非有限 norm 应重抽，连续 10 次失败则报错。

\subsection{语义 queries 与增广表}
持久参数为 $Q^U,Q^F\in\R^{K\times d}$ 和 $q^C\in\R^d$，$K=32$。
为每行增加 $K$ 个 Unit Query 列，为每列增加 $K$ 个 Feature Query 行：
\begin{equation}
 H^{(0)}\in\R^{(N+K)\times(M+K)\times d},\qquad
 h_i^{(0)}=\begin{cases}
 h^{\rm val}_{ra},&i=(r,a)\in\V,\\
 q^C,&i=(r,a)\in\Q,\\
 Q^U_{k,:},&i=(r,M+k),\ r\le N,\\
 Q^F_{k,:},&i=(N+k,a),\ a\le M,\\
 0,&\text{其他地址。}
 \end{cases}
 \label{eq:compiler}
\end{equation}
所有非事实 queries \textbf{只接收、不供源}。原始缺失且未被查询的位置，以及右下 $K\times K$ 角都是 Null。
一个 seed 的第 $k$ 槽跨所有行或列共享，没有每个数据集专属的持久 bank。

\begin{figure}[htbp]\centering
\begin{tikzpicture}
\node[fact,minimum width=78mm,minimum height=31mm] (cells) {原始 Cell 平面 $N\times M$\\Value / Query / Null};
\node[choice,minimum width=48mm,minimum height=31mm,right=0mm of cells] (u) {Unit Queries\\$N\times K$};
\node[choice,minimum width=78mm,minimum height=20mm,below=0mm of cells] (f) {Feature Queries $K\times M$};
\node[box,minimum width=48mm,minimum height=20mm,right=0mm of f] (n) {Null $K\times K$};
\node[font=\small,align=center,below=6mm of $(f.south)!0.5!(n.south)$]
 {每个位置存 $d=384$ 维 carrier；$K=32$ 是语义 bank 槽数。\\
 $S=128$ inducing 摘要位于 block 内部，不是这张表的新行或新列。};
\end{tikzpicture}
\caption{表格地址与语义 banks。Null 始终保持向量零；数值为零的真实观测仍是 Value。}
\label{fig:carrier}
\end{figure}

\section{一个完整的 OAttention block}
\label{sec:block}

本节精确定义所有子层。取 receiver $h_i\in\R^d$，候选 source $x_j\in\R^d$，
资格 mask $m_j\in\{0,1\}$。原始表的 $m_j$ 由固定 $\V$ 决定；摘要回读的 mask 由本层列证据状态决定。

\subsection{Presence、源集合和多头注意力}
默认 presence 函数与实际 source 集合为
\begin{equation}
 \rho(h)=\frac{\|h\|_2^2}{1+\|h\|_2^2},\qquad
 \mathcal A=\{j:m_j=1,\ \rho(x_j)>0\}.
 \label{eq:presence}
\end{equation}
资格 mask 不能用 $\rho$ 代替。非零 Query 仍不允许供源。
$\rho$ 读取\emph{投影前} carrier，默认 $\tau_{\rm pres}=1$，不是 norm 与投影后的量。

每个 OMAB 有独立的四个仿射投影 $P_Q,P_K,P_V,P_O$，各含 $d\times d$ 权重与 $d$ bias。
合并投影输出分成 8 个 head，每头 $d_h=48$：
\begin{align}
 q_i^{(h)}&=[P_Q(h_i)]_h,&k_j^{(h)}&=[P_K(x_j)]_h,&v_j^{(h)}&=[P_V(x_j)]_h,\\
 u_{ij}^{(h)}&=\frac{\langle q_i^{(h)},k_j^{(h)}\rangle}{\sqrt{48}}+\log\rho(x_j),
 &&j\in\mathcal A.\label{eq:attn-logit}
\end{align}
默认直接投影 carrier，\textbf{不另加 Q/K pre-norm、cosine attention、RoPE 或 QASSMax}。
这些选项会改变函数，不是照搬参考模型时可以隐含加入的细节。

若 $\mathcal A\ne\varnothing$，只在此集合计算稳定 softmax：
\begin{equation}
 a_i=\rho(h_i)\,P_O\!\left(\operatorname{Concat}_{h=1}^{8}
       \sum_{j\in\mathcal A}\sm_{j}(u_{ij}^{(h)})v_j^{(h)}\right).
 \label{eq:attention}
\end{equation}
若 $\mathcal A=\varnothing$，直接置 $a_i=0$，不执行空 softmax，也不保留 $P_O$ 的 bias。
非空时 receiver gate 位于整个 output projection 之后。
源质量只通过 $\log\rho$ 进入权重，不再次乘到 $v_j$ 上。

\subsection{Token-local residual 与 FFN}
每个 OMAB 的唯一归一化位于 FFN 前，有 gain $\gamma\in\R^d$，无可训练 shift：
\begin{align}
 r_i&=h_i+a_i,\\
 \Norm(r_i)&=\gamma\odot\frac{r_i}{\sqrt{d^{-1}\|r_i\|_2^2+10^{-6}}},\\
 \FFN(t)&=B_2\,\operatorname{GELU}(B_1t+b_1)+b_2,\\
 \OMAB(h_i;X)&=\begin{cases}
 r_i+\rho(r_i)\FFN(\Norm(r_i)),&i\notin\Nul,\\
 0,&i\in\Nul.
 \end{cases}\label{eq:omab}
\end{align}
$B_1\in\R^{768\times384}$，$B_2\in\R^{384\times768}$，bias 分别为 768、384 维。
GELU 取精确式 $t\Phi(t)$，dropout 为零；默认没有 block 后额外 normalization。
Attention receiver gate 读取 $h_i$，FFN gate 读取 $r_i$，两者不可混用。
空 source 时 attention 为零，但非 Null receiver 的 FFN 仍执行。

\begin{figure}[htbp]\centering
\begin{tikzpicture}[node distance=7mm]
\node[box,text width=28mm] (h) {receiver $h$};
\node[box,text width=34mm,right=of h] (a) {OAttention\\源 mask + presence};
\node[box,text width=22mm,right=of a] (r) {$r=h+a$};
\node[box,text width=35mm,below=11mm of a] (ff) {RMSNorm $\to$ FFN\\乘 $\rho(r)$};
\node[fact,text width=33mm,right=of ff] (o) {residual 相加\\Null 精确写零};
\node[fact,text width=28mm,above=7mm of a] (s) {合法 sources $X$};
\draw[arr] (h)--(a);\draw[arr] (s)--(a);\draw[arr] (a)--(r);
\draw[arr] (r.south)--++(0,-4mm)-|(ff.north);\draw[arr] (ff)--(o);
\draw[arr,dashed] (r.south)--(o.north);
\end{tikzpicture}
\caption{一个 OMAB 的两条 residual 更新。图中 $r$ 同时作为最终 residual 的基底；
Norm 只沿当前 token 的 384 维归约。所有 projection biases 都包含在数学定义与参数统计中。}
\label{fig:omab}
\end{figure}

\subsection{默认：列内固定 inducing，随后直接行内交互}
第 $\ell$ 层有独立 learned seed $I^{(\ell)}\in\R^{128\times384}$，同层所有列共享。
令 $X_a^{(\ell)}=(h_{sa}^{(\ell)})_{s\in\mathcal C_a}$，
$e_a^{(\ell)}=\ind[\exists s\in\mathcal C_a:\rho(h_{sa}^{(\ell)})>0]$。
先完成全表列更新，再完成全表行更新：
\begin{align}
 G_a^{(\ell)}&=\begin{cases}
 \OMAB_{\ell,c}(I^{(\ell)};X_a^{(\ell)}),&e_a^{(\ell)}=1,\\
 0_{S\times d},&e_a^{(\ell)}=0,
 \end{cases}\label{eq:collect}\\
 H^{(\ell,\rm col)}_{:,a}&=\OMAB_{\ell,d}(H^{(\ell)}_{:,a};G_a^{(\ell)}),\label{eq:read}\\
 H^{(\ell+1)}_{r,:}&=\OMAB_{\ell,r}
 \left(H^{(\ell,\rm col)}_{r,:};(H^{(\ell,\rm col)}_{r,b})_{b:(r,b)\in\V}\right).
 \label{eq:row}
\end{align}
式~\eqref{eq:collect} 的 seeds 都是合法非 Null receivers；空证据列必须覆盖整个 collect 结果为零，
不能让 seed residual 冒充证据。回读仅使用非零摘要的 presence；摘要不进入 $\V$ 或 terminal 支持集合。
在右侧 Unit Query 列，原始事实为空，$G=0$；底部 Feature Query 行也没有 row sources。
这些 receivers 仍可在没有跨地址信息的子层中执行局部 FFN。

同层 collect、read、row 的全部权重相互独立，层间也独立；每个子层跨地址共享权重。
$G$ 是当前 episode 的 activation，不能保存成 learned parameter。每一层从自己的 $I^{(\ell)}$ 开始收集，
不把上一 episode 或上一层的 $G$ 当成初始 seed。
\textbf{即使 $N<128$ 也使用全部 128 个槽位}，不取 $\min(N,S)$，不自动切换 direct attention。

\subsection{保留 without inducing 的两种公平对照}
Direct baseline 的每层为
\begin{align}
 H^{(\ell,\rm col)}_{:,a}&=\OMAB_{\ell,\rm col}
       (H^{(\ell)}_{:,a};X_a^{(\ell)}),\\
 H^{(\ell+1)}_{r,:}&=\OMAB_{\ell,\rm row}
       (H^{(\ell,\rm col)}_{r,:};(H^{(\ell,\rm col)}_{r,b})_{b:(r,b)\in\V}).
 \label{eq:direct}
\end{align}
它没有 inducing seed 参数，语义 banks 和 typed terminal 保留。不是设 $S=0$，也不是从一个
with-inducing checkpoint 临时跳过两个子层。

\begin{figure}[htbp]\centering
\begin{tikzpicture}[node distance=6mm]
\node[font=\bfseries\small,anchor=west] at (-1,1) {Without inducing：直接轴向 baseline};
\node[fact,text width=25mm] (x) {$H^{(\ell)}$};
\node[box,text width=42mm,right=of x] (dc) {直接 column OMAB\\同列 visible sources};
\node[box,text width=42mm,right=of dc] (dr) {直接 row OMAB\\同行 visible sources};
\draw[arr] (x)--(dc);\draw[arr] (dc)--(dr);
\node[font=\bfseries\small,anchor=west] at (-1,-1.4) {With inducing：当前默认};
\node[fact,text width=25mm,below=22mm of x] (xx) {$H^{(\ell)}$};
\node[choice,text width=27mm,right=of xx] (cc) {Collect\\$X_a\to G_a$};
\node[choice,text width=27mm,right=of cc] (dd) {Read\\$G_a\to H_{:,a}$};
\node[box,text width=27mm,right=of dd] (rr) {Row OMAB\\直接行交互};
\draw[arr] (xx)--(cc);\draw[arr] (cc)--(dd);\draw[arr] (dd)--(rr);
\node[font=\small,below=6mm of $(xx.south)!0.5!(rr.south)$,align=center]
 {$S=128$ 是通信瓶颈；$K=32$ 的语义 banks 在两条路径中都存在。};
\end{tikzpicture}
\caption{从直接轴向模型到固定槽位模型。下图的 read 同时读取原来的 receiver residual；
图展示通信路径，不表示串行丢弃原 $H$。两种结构需要分别构建与训练。}
\label{fig:withwithout}
\end{figure}

同深度 direct 取 $L=15$，比较结构代价；近似同参数量 direct 取 $L=23$，比较参数预算。
后者同时增加路径深度，不能宣称是仅改变 inducing 的单因素实验。两条比较应分别报告。

\section{Typed additive readout 与预测}
\label{sec:readout}

\subsection{从地址读出三个视图}
从最后一层按地址读取，不增加 token 交互：
\begin{equation}
 c_{ra}=H^{(L)}_{r,a,:}\in\R^d,\quad
 U_r=H^{(L)}_{r,M+1:M+K,:}\in\R^{K\times d},\quad
 F_a=H^{(L)}_{N+1:N+K,a,:}\in\R^{K\times d}.
 \label{eq:views}
\end{equation}
持久 readout 参数为 $W\in\R^{K\times d}$、$w_F,w_U\in\R^d$，无 bias：
\begin{equation}
 \boxed{z_{ra}=Wc_{ra}+\lambda_F F_aw_F+\lambda_U U_rw_U\in\R^{32}.}
 \label{eq:response}
\end{equation}
本文完整配置取固定 gates $\lambda_F=\lambda_U=1$。它们独立，不训练、不作 schedule；
两个 banks 是当前 forward 的表示，不是作用于 $c$ 的动态权重。
同一公式用于 Query 与 visible support；不能只为 Query 计算 $z$。

\subsection{匹配：同列 Gaussian 权重}
对目标 $t=(r,a)\in\Q$ 及 $s\in\mathcal C_a$，定义
\begin{equation}
 \delta_s=z_{sa}-z_{ra},\quad
 \psi_s=-\frac{\|\delta_s\|_2^2}{2\tau_{\rm match}^2},\quad
 \alpha_s=\sm_{s\in\mathcal C_a}(\psi_s),\qquad\tau_{\rm match}=1.
 \label{eq:matching}
\end{equation}
重复样本按不同地址保留，不能去重；只沿同列合法支持归一化，不跨列、不含 Query 或 Null。
带宽固定、跨列共享；模型能够改变响应尺度，所以固定带宽并不消除尺度与 conditioning 的训练问题。
默认不做 top-$k$ 检索、不 detach supports，也不以 attention learned $Q/K$ 替代此式。

\subsection{Numeric：可微的带截距局部线性拟合}
用可见值定义
\begin{equation}
 \mu_a=\frac1{n_a}\sum_s x_{sa},\quad
 \sigma_a=\sqrt{10^{-12}+\frac1{n_a}\sum_s(x_{sa}-\mu_a)^2},\quad
 y_s=\frac{x_{sa}-\mu_a}{\sigma_a}.
 \label{eq:target-scale}
\end{equation}
这里的 scale 与 encoder 的 $\widehat\sigma_a+10^{-6}$ 不相同。
取 $\lambda_{\rm LL}=10^{-2}$，每个 Query 独立求解
\begin{equation}
 (\eta,v)=\arg\min_{\eta\in\R,v\in\R^{32}}
 \sum_s\alpha_s(y_s-\eta-v^\top\delta_s)^2+10^{-2}\|v\|_2^2,
 \qquad\widehat x_{ra}=\mu_a+\sigma_a\eta.
 \label{eq:ll}
\end{equation}
截距不惩罚。局部 $(\eta,v)$ 是 activation，不是跨 episode 训练的参数。
实现不求显式矩阵逆，而使用消去截距的形式：
\begin{align}
 \bar\delta&=\sum_s\alpha_s\delta_s,&\bar y&=\sum_s\alpha_sy_s,\\
 e_s&=\delta_s-\bar\delta,&\widetilde y_s&=y_s-\bar y,\\
 C&=\sum_s\alpha_s e_se_s^\top,&g&=\sum_s\alpha_s e_s\widetilde y_s,\\
 v&=\operatorname{CholeskySolve}(C+10^{-2}I,g),&
 \eta&=\bar y-v^\top\bar\delta.
 \label{eq:solve}
\end{align}
FP64 作为首个 correctness 实现的 terminal 精度；参数主干可用 FP32。
拟合矩阵对称化为 $(C+C^\top)/2$ 以移除求和舍入产生的非对称。
若 Cholesky 失败，应报告失败及输入统计，不悄悄增大 ridge、切换伪逆或返回 NW。
这些替代都会改变已经定义的模型。

\begin{proposition}[局部解唯一]
只要支持权重非负、和为 1 且输入有限，正 ridge 使式~\eqref{eq:ll} 的解唯一。
这不要求支持数大于 32。
\end{proposition}
\begin{proof}
$C$ 是加权协方差，故 $C\succeq0$，$C+10^{-2}I\succ0$，斜率有唯一解；
截距由 $\eta=\bar y-v^\top\bar\delta$ 唯一确定。
这也解释了 $n_a\ge2$ 是本模型的证据策略，而不是矩阵可逆的最低条件。
\end{proof}

LL 可以外推到支持值域之外；本稿只承诺 numeric 点预测，不定义分位数或不确定性。
若需要严格正值、整数等输出，应另行设计 domain-safe terminal，不能无声明地裁剪。

\subsection{Nominal 与 ordinal：完整声明域的平滑 NW}
两类离散列共用概率公式，ordinal 的顺序目前只用于输入编码：
\begin{equation}
 m(b)=\sum_s\alpha_s\ind[x_{sa}=b],\qquad
 \widehat p_{ra}(b)=\frac{m(b)+\varepsilon_{\rm cat}}{1+D_a\varepsilon_{\rm cat}},
 \quad b\in\mathcal D_a,\quad\varepsilon_{\rm cat}=10^{-6}.
 \label{eq:pmf}
\end{equation}
概率和为 1，各坐标为正；默认点分类为按 schema 顺序打破并列的 argmax。
未见合法类别获得 $\varepsilon_{\rm cat}/(1+D_a\varepsilon_{\rm cat})$，不能据此声称具有 zero-shot 类别识别能力。
参数量不随类别数增加，但 PMF 的存储与输出时间仍随 $D_a$ 增长。
未设置训练外类别数能力上限；初始训练分布的类别数范围在后文单独指定。

\begin{figure}[htbp]\centering
\begin{tikzpicture}[node distance=7mm and 8mm]
\node[box,text width=34mm] (c) {Cell $c_{ra}$\\$Wc_{ra}$};
\node[box,text width=34mm,below=of c] (f) {Feature $F_a$\\$F_aw_F$};
\node[box,text width=34mm,below=of f] (u) {Unit $U_r$\\$U_rw_U$};
\node[choice,text width=26mm,right=of f] (z) {相加\\$z_{ra}\in\R^{32}$};
\node[box,text width=38mm,right=of z] (m) {同列 $z_{sa}-z_{ra}$\\Gaussian 权重 $\alpha_s$};
\node[fact,text width=38mm,above=of m] (n) {Numeric\\weighted LL 点预测};
\node[fact,text width=38mm,below=of m] (d) {Nominal / ordinal\\平滑 NW 概率};
\draw[arr] (c.east)--(z.west);\draw[arr] (f)--(z);\draw[arr] (u.east)--(z.west);
\draw[arr] (z)--(m);\draw[arr] (m)--(n);\draw[arr] (m)--(d);
\end{tikzpicture}
\caption{响应与 terminal 分为两个接口。支持的 $z_{sa}$ 使用完全相同的响应公式；
LL/NW 的被拟合值是同列原始可见值，不是重新编码后的 value 向量。}
\label{fig:readout}
\end{figure}

\subsection{必须面对的取舍：Feature 公共平移不影响当前预测}
同列 Query 与 supports 使用同一个 $F_a$，因此
\begin{equation}
 \delta_s=W(c_{sa}-c_{ra})+\lambda_U(U_s-U_r)w_U.
 \label{eq:translation}
\end{equation}
Feature 项精确抵消。Gaussian 权重和中心化 LL 都只读取这些差值；NW 则只读取权重与可见标签。
故在\emph{固定其他参数、证据与 codebook} 时，$\lambda_F=0$ 与 1 产生相同预测。
$\lambda_U$ 不具有一般的同列平移抵消性质。

\begin{decision}
本稿保留 $\lambda_F=\lambda_U=1$ 和完整 Feature 接口，以忠实表达 TAR 当前共同响应结构；
但不声称三个分支都已从此损失学到有效功能。默认预测损失对 $Q^F,w_F$ 的梯度为零。
相应 12,672 个 Feature 专属参数占总量约 0.0234\%。
可以在实现优化中省略其计算，但必须将完整结构作为对照，并验证预测与其余梯度等价。
\end{decision}

原因不仅是加法：Feature Queries 没有供源资格，不能借其他路径影响 Cell/Unit；
它们对共享 block 参数的贡献也在此 terminal 下抵消。共享参数仍由 Cell/Unit 路径训练。
优化器不能通过 weight decay 制造“Feature 已被监督学习”的假象；本稿不衰减任何 semantic seeds 或 readout vectors。

\texttt{explorations/kernel-attention-readout.tex} 探索的 learned matching 可能改变这一结论。
那是下一项可单独检验的设计，不在本稿默认中偷偷启用。若未来采用它，必须同时重算参数、梯度、
对称性和 terminal 行为，不能只替换一个名字。

\section{数值默认、初始化与参数预算}
\label{sec:init}

表~\ref{tab:defaults} 将影响模型函数的标量全部固定。它们是本稿的可修订起始配置。
\begin{table}[htbp]\centering\small
\caption{首份配置的数值默认。}\label{tab:defaults}
\begin{tabularx}{\linewidth}{@{}l l Y@{}}\toprule
设置 & 数值 & 作用 \\\midrule
$d,L,h,d_{\rm ff}$ & $384,15,8,768$ & 同层 collect/read/row 独立参数 \\
$K,d_z,J$ & $32,32,96$ & 语义 bank、响应宽度、Fourier 频率数 \\
$S$ & 128 & 每层固定，跨列共享 seeds \\
$\lambda_F,\lambda_U$ & $1,1$ & 两个固定独立 gates \\
$\tau_{\rm pres}$ & 1 & $\rho(h)=\|h\|^2/(1+\|h\|^2)$ \\
$\varepsilon_{\rm norm}$ & $10^{-6}$ & RMS denominator \\
$\varepsilon_{\rm scale},\sigma_{\min}$ & $10^{-6},10^{-6}$ & encoder 与 terminal 不同公式 \\
$\tau_{\rm match}$ & 1 & Gaussian 带宽 \\
$\lambda_{\rm LL}$ & $10^{-2}$ & 只惩罚 LL slope \\
$\varepsilon_{\rm cat}$ & $10^{-6}$ & 每个声明类别的平滑量 \\
Dropout / estimators & $0 / 1$ & 单次确定性 forward，固定 codebook \\
额外位置或长度机制 & 关闭 & 不加 RoPE、QASSMax、feature grouping \\\bottomrule
\end{tabularx}
\end{table}

\subsection{从初始化到第一步反向传播}
定义 $B=3L=45$ 为默认 OMAB 个数，direct 用 $B=2L$。
全局初始化 RNG seed 为 20260906；不同模块用独立子流，不同 baseline 单独初始化。
\begin{enumerate}
\item 所有 $P_Q,P_K,P_V$ 权重用 Xavier uniform，所有 bias 为零。
$P_O$ 与 FFN 第二层先 Xavier uniform，再乘 $1/\sqrt{2B}$；FFN 第一层用 Xavier uniform。
这是 residual 初始尺度选择，不是在每次 forward 额外乘一个 gate。
\item RMS gain 初始化为 1。无 trainable shift。$Q^U,Q^F,q^C,I^{(\ell)}$ 每个向量独立抽标准 Gaussian 后归一化到 norm 1。
此后直接作为无约束可学习向量，不每步重新归一化；初始 $\rho=1/2$，避免零 seed 无法激活。
\item 频率 $\omega_j=2^{-3+6(j-1)/(J-1)}$，直接训练这 $J$ 个实数。
$A$ 用 Gaussian 矩阵的 thin QR 得到正交列，再除以 $\sqrt J$。
因此初始化时每个 continuous carrier 的 norm 为 1。QR 的列符号固定为 $R$ 对角线非负，便于重放。
\item $W,w_F,w_U$ 的每个分量独立取 $\mathcal N(0,1/(Kd))$，不初始化为零。
同一坐标尺度便于起步，但不保证训练后匹配仍处于良好条件区间。
\end{enumerate}

这些初值不构成训练稳定性证明。每次编译检查所有非 Null 初始 carrier 有限且非零；
训练中出现精确零 Value 应记录并停止诊断，不用任意 presence 阈值把它隐藏。
不要将范数很小误判为精确零；稳定 norm、log-presence 与 FP32/64 累加应先保证数值实现正确。

\subsection{逐张量参数量，而不是按名字估计}
一个 OMAB 有四个带 bias 的 $d\to d$ 投影、一个仅 gain 的 RMSNorm、两个带 bias 的 FFN 线性层：
\begin{equation}
 P_{\rm OMAB}=4d^2+2dd_{\rm ff}+6d+d_{\rm ff}=1,182,720.
 \label{eq:count-omab}
\end{equation}
共享 encoder、semantic seeds 和 readout 合计
\begin{equation}
 P_{\rm shared}=J+2Jd+(2K+1)d+(K+2)d=111,840.
 \label{eq:count-shared}
\end{equation}
故默认模型
\begin{equation}
 P=L(3P_{\rm OMAB}+Sd)+P_{\rm shared}
 =\boxed{54,071,520}.
 \label{eq:parameter-count}
\end{equation}
Norm、bias、Fourier 频率、inducing 与语义 seeds 全部计入；codebook、$H,G,F,U$、局部 LL coefficients
及固定超参数不计入。总量比报告舍入的 53M 高约 2.0\%，比 58M 低约 6.8\%。
这是标称参数匹配，不是 checkpoint 实测值的精确对齐。

\begin{table}[htbp]\centering\small
\caption{独立构建的三个配置。}\label{tab:budgets}
\begin{tabularx}{\linewidth}{@{}l r r Y@{}}\toprule
配置 & $L$ & 参数 & 比较用途 \\\midrule
With inducing 默认 & 15 & 54,071,520 & 首份实现目标 \\
Direct 同深度 & 15 & 35,593,440 & 槽位替换的结构比较 \\
Direct 近似同参数 & 23 & 54,516,960 & 参数接近，深度同时变化 \\\bottomrule
\end{tabularx}
\end{table}

运行 \texttt{python3 design/parameter\_count.py} 会独立枚举 608 个默认参数张量，
再与上述闭式公式核对。它不加载模型、不声称代码已经实现。
Feature 专属参数仍计入 declared trainable tensors；其默认 loss 梯度为零已在前节说明。

\section{训练目标与可执行的起始配方}
\label{sec:training}

\subsection{唯一的监督入口}
每个训练 Query 的损失为
\begin{equation}
 \ell_{ra}=\begin{cases}
 \frac12[(\widehat x_{ra}-x^\star_{ra})/\sigma_a]^2,&a\text{ numeric},\\
 -\log\widehat p_{ra}(x^\star_{ra}),&a\text{ nominal 或 ordinal}.
 \end{cases}
\end{equation}
对非空类型分支分别平均，再相加；episode 之间等权平均：
\begin{equation}
 \mathcal L(E)=\sum_{b\in\{\rm num,disc\}:|\Q_b|>0}
       \frac1{|\Q_b|}\sum_{(r,a)\in\Q_b}\ell_{ra},
 \qquad \mathcal J(\theta)=\mathbb E_{E,\xi_{\rm prep}}\mathcal L(E).
 \label{eq:loss}
\end{equation}
不是把所有 Cell 混成一个均值，也不是把两类分支再次除以 2。
LL 的 $\alpha,C,g,v,\eta$ 全程保留 autograd；可见统计和 codebook 是常量。
Query 与 support 的 responses 都有梯度。若 $Av=g$，微分为
$dv=A^{-1}(dg-dA\,v)$；禁止 detach $A$ 或把求解当作不训练的 sklearn 后处理。

\subsection{Mask sampler 的完整默认}
对给定自然观测地址 $\mathcal O$，生成只依赖地址的均匀随机排列。
目标数量为
\begin{equation}
 q_{\rm goal}=\min(64,\max(1,\lfloor0.15|\mathcal O|\rfloor)).
\end{equation}
按排列扫描：仅当隐藏该 Cell 后所在列仍至少留两条 visible observations 时，才把它加入 $\Q$；
达到 $q_{\rm goal}$ 或扫描结束即停止。$\Q$ 为空则返回 \texttt{no-valid-episode}，重新抽表。
最终 $\V=\mathcal O\setminus\Q$；保留未达到目标数量的合法非空 episode。
该算法完全确定了 sampler，而没有声称它在所有合法 Query 子集中均匀。

自然缺失与人为 mask 分开保存。预处理在 mask 后才拟合；每次 episode 重抽 nominal codes，
不把隐藏类别加入 codebook。默认只有 masked-cell 训练，不暗中加入监督目标列训练或对比学习辅助损失。
监督 ICL 转移是否需要结构化 mask，是后续实验问题。

\subsection{一个可从零生成的训练表分布}
\label{sec:prior}
为了让实现不依赖尚未选定的数据文件，本文规定一个\textbf{最小 synthetic prior 起始配方}。
它服务于首个训练循环与数学诊断，不声称复现 TabPFN-3 的预训练 prior，也不声称它足以训练成功的基座模型。
未来正式预训练可能需要更丰富的表分布及训练规模；模型前向定义不依赖这一选择。

每个 episode 的采样顺序如下；所有 Uniform 离散集合默认等概率，所有连续区间含端点与否不影响实数分布：
\begin{enumerate}
\item 独立抽 $N\in\{64,128,256,512\}$、$M\in\{4,8,16,32\}$。
创建 $H_0=8$ 个隐藏节点和 $M$ 个输出节点，按 $j=1,\ldots,H_0+M$ 的拓扑序生成。
\item 节点 $j$ 的父节点数 $p_j$ 从 $\{0,\ldots,\min(3,j-1)\}$ 均匀抽取；
从前 $j-1$ 个节点无放回选父集合。$w_{ji}\sim\mathcal N(0,1/\max(1,p_j))$，
$b_j\sim U[-0.5,0.5]$，$a_j\sim\operatorname{LogUniform}[0.5,2]$，
$\nu_j\sim\operatorname{LogUniform}[0.05,0.5]$，函数 $f_j$ 从 identity、tanh、sin、ReLU 均匀选择。
\item 对每行独立生成 $\epsilon_{rj}\sim\mathcal N(0,1)$。根节点取 $z_{rj}=\epsilon_{rj}$；
非根节点取
\begin{equation}
 z_{rj}=a_j f_j\!\left(b_j+\sum_{i\in\operatorname{pa}(j)}w_{ji}z_{ri}\right)+\nu_j\epsilon_{rj}.
 \label{eq:synthetic}
\end{equation}
表值使用最后 $M$ 个节点，前 8 个节点不进入输入。根节点未用到的随机参数不影响生成。
\item 每个输出列独立选择 numeric/nominal/ordinal，概率为 $0.5/0.3/0.2$。
Numeric 取 $x_{ra}=s_a z_{r,H_0+a}+m_a$，$s_a\sim\operatorname{LogUniform}[0.1,10]$、$m_a\sim U[-2,2]$。
离散列从 $\{2,4,8,16\}$ 选择类别数 $D_a$，用预先固定的标准正态分位点
$\Phi^{-1}(k/D_a)$，$k=1,\ldots,D_a-1$，把 latent 值分桶。
Ordinal 保留桶次序；nominal 施加独立随机类别重标。声明域在抽样列值之前由 $D_a$ 给定。
\item 为每列抽自然缺失概率 $p_a\sim U[0,0.2]$，各 Cell 独立缺失；
再独立打乱输出行列及其 schema/ID，执行上述 mask sampler。
分类域允许有本 episode 未出现的类别，不根据 Query 真值重新生成分桶或类别域。
\end{enumerate}

节点机制参数在同一表的行间共享，每张表重抽；噪声在行间独立。
训练期 RNG 使用互相独立的 prior、missingness、mask、codebook 子流。
例如用 NumPy \texttt{SeedSequence([20260906, split\_id, episode\_id])} 派生这些子流并保存参数，
分布复现以保存生成的 episode/参数为准。验证使用独立 split ID，不能只换 mask 而复用同一潜在生成世界。
这里的 SCM 仅用于产生合成相关结构，不给 TAR 内部响应赋予因果识别语义。

\subsection{优化器与初始训练预算}
下列数值是\textbf{待验证的运行起点}，不是已完成实验或已获批准的算力作业：
\begin{table}[htbp]\centering\small
\begin{tabularx}{\linewidth}{@{}l l Y@{}}\toprule
选项 & 默认 & 精确定义 \\\midrule
Optimizer & AdamW & $\beta=(0.9,0.95)$，$\epsilon=10^{-8}$ \\
Learning rate & $10^{-4}$ & 前 2,000 optimizer steps 线性 warm-up \\
总更新数 & 100,000 & warm-up 后 cosine 到 $10^{-5}$ \\
有效 batch & 16 episodes & microbatch 1，累积 16 次，每次 loss 除以 16 \\
Weight decay & $10^{-2}$ & 仅 $A,W$、attention/FFN 权重矩阵 \\
不衰减 & 0 & 所有 bias、norm gain、频率、seeds、$w_F,w_U$ \\
Gradient clip & 1 & 所有累积梯度 global norm，一次 optimizer step 前执行 \\
训练精度 & FP32 & terminal FP64；BF16 是验收后的工程选项 \\
模型随机性 & dropout 0 & 每个 episode 新 codebook；无 augmentation ensemble \\\bottomrule
\end{tabularx}
\end{table}

用 $t=1,\ldots,T$ 编号 optimizer steps，$T=100000$、$w=2000$、$\eta_{\max}=10^{-4}$、
$\eta_{\min}=10^{-5}$：
\begin{equation}
 \eta_t=\begin{cases}
 \eta_{\max}t/w,&t\le w,\\
 \eta_{\min}+\tfrac12(\eta_{\max}-\eta_{\min})
 [1+\cos(\pi(t-w)/(T-w))],&t>w.
 \end{cases}
\end{equation}
DDP 时按全局 episode 数 16 配置累积和 loss scaling，不能额外乘一次 world size。
若资源要求调整 batch，需同时记录新设置，不把未达到的有效 batch 写入配置。
每次更新若出现非有限 loss/gradient，应停止并保存失败样例；不静默跳步后仍称完成原预算。

初步验证固定 256 个独立 synthetic episodes，每 1,000 步报告 numeric/discrete losses 和失败数量，
不以训练 best single episode 代替泛化。先做第~\ref{sec:acceptance} 节小规模验收，再决定是否启动上述预算。
此文件不启动训练、不修改模型 registry，也不产生能力证明。

\section{推理、缓存与资源边界}
\label{sec:inference}

\subsection{最小正确推理路径}
固定模型参数与 codebook，按式~\eqref{eq:input} 验证输入并编译整表；
逐层完成 collect、read、row；只在全部层完成后读出 Query 和支持 responses；
最后按目标列类型调用 terminal。推理不对当前表进行梯度更新。
可以只物化需要的最终 $z$，但不能漏算形成这些 $z$ 所依赖的 Value carriers。

默认一个 estimator。分类返回完整声明域 PMF 和预测标签，回归返回点值；
输出同时包含 target 地址、状态、支持数。调试可返回 $z$、matching entropy、有效支持数
$n_{\rm eff}=1/\sum_s\alpha_s^2$ 与 LL 系统条件数，但它们不是额外监督输入。

\subsection{分块必须保持全局证据归一化}
初始 correctness 实现保留完整 $H^{(\ell)}$；默认 receiver row chunk 为 256、collect column chunk 为 4、
terminal query chunk 为 32、source chunk 为 2048。这些是 TAR 的起始内存参数，不改变数学超参数。
不足一个 chunk 的尾部按实际长度处理，padding 不进入源集合。

Collect 可沿列分块，因为不同列独立。沿 sources 分块时必须合并\emph{同一个 softmax} 的统计，
不能先各块归一化再平均。每个 head/query 保存最大 logit $m$、分母 $Z$ 和加权值和 $U$。
两块 $(m_1,Z_1,U_1)$、$(m_2,Z_2,U_2)$ 合并为
\begin{equation}
 m=\max(m_1,m_2),\quad
 Z=e^{m_1-m}Z_1+e^{m_2-m}Z_2,\quad
 U=e^{m_1-m}U_1+e^{m_2-m}U_2.
 \label{eq:stream}
\end{equation}
全部合法 sources 合并后才计算 $U/Z$、output projection、receiver gate 和 FFN；
空块作为恒等空统计跳过，全部为空则执行零更新分支。
Terminal 的 Gaussian softmax 同样需要全支持归一化。为 LL 计算中心协方差时，首个实现用 FP64 两遍归约，
避免原始二阶矩相减造成消减误差；优化版需与此参考核对。

\begin{figure}[htbp]\centering
\begin{tikzpicture}[node distance=7mm]
\node[fact,text width=37mm] (old) {本层全部 visible carriers\\$H^{(\ell)}_{\V}$};
\node[choice,text width=37mm,right=of old] (g) {全证据 Collect\\$G^{(\ell)}_a$};
\node[box,text width=37mm,right=of g] (stream) {按行 Read + Row\\生成 $H^{(\ell+1)}$};
\draw[arr] (old)--(g);\draw[arr] (g)--(stream);
\node[box,text width=119mm,below=11mm of g] (barrier)
 {层间同步：下一层必须收集本层更新后的全部 visible carriers。\\
 不能固定第一层摘要，把一个 row chunk 独自跑完 15 层。};
\draw[arr] (stream.south)|-(barrier.east);
\draw[arr] (barrier.west)-|([xshift=-5mm]old.west)--(old.west);
\end{tikzpicture}
\caption{TAR 的逐层全局摘要与分块回读。交替的行更新会改变下一层列证据，
所以不能直接套用另一架构的整网两阶段缓存结论。}
\label{fig:cache}
\end{figure}

\subsection{缓存何时有效}
缓存 key 必须包含模型参数、schema、全部可见地址和值、预处理 realization、层状态及数值配置。
仅在这些不变时，已有 Value/Unit/Feature carriers 与 summaries 才可复用。
向原来 Null 的位置新增 Query 且不改变 $\V$，只需计算新增 Query 的 receiver 路径；
已有 queries 不供源，所以不能改写已有响应。批量 Query 拆分需要固定完整可见事实和 codebook。

修改 visible value、加入测试行 covariates、改变 masking、重新抽 codebook 或更新权重都会使相关缓存失效。
因此第~\ref{sec:input} 节的逐测试行 API 默认不能共享一份“训练表已完成全部层”的 KV cache。
第一份实现宁可重新前向；若要引入训练表独立缓存，必须另外设计信息边界并证明对应的新合同。

\subsection{参数相近不意味着显存或 FLOPs 相近}
令 $\widetilde N=N+K$、$\widetilde M=M+K$、$T=\widetilde N\widetilde M$。
计入线性投影和 FFN，每层 dense reference 的量级为
\begin{align}
 C_{\rm induced}&=O\big((T+\widetilde MS)(d^2+dd_{\rm ff})
                    +T(S+\widetilde M)d\big),\\
 C_{\rm direct}&=O\big(T(d^2+dd_{\rm ff})+T(\widetilde N+\widetilde M)d\big).
 \label{eq:cost}
\end{align}
这是增广 dense 上界；实际 sparse kernels 可跳过不供源的地址。
固定 $M,K,S,d$ 时 inducing backbone 关于 $N$ 的 attention 项为线性，但并未消除关于特征数的二次 row attention。
$N<S$ 时 collect/read 的常数成本可能大于 direct；保留固定槽是模型选择，不是加速保证。

单个 $H$ 就含 $Td$ 个数。默认 $N=512,M=32$ 时，FP32 单份 $H$ 约 51 MiB；
15 层训练还需要中间量、反向存储、优化器与工作区。
仅 FP32 参数、梯度和 Adam 两组矩已约 825 MiB，不含 activations。
若 $N=10^6,M=200$，单份 BF16 dense $H$ 已约 166 GiB。
所以本稿\textbf{没有百万行单卡运行承诺}。分块 score 可以降低注意力工作区，却不会自动消除保存完整 $H$ 的成本。

可设计逐层 CPU/offload 或重计算，使 accelerator 上只保留当前 chunk 和本层 summaries；
这仍需要整体状态存储和每层全证据扫描，并产生 I/O 或重算成本。
当前文档给出等价归约规则，尚未实现这类高性能后端。

对 $Q=|\Q|$ 个目标、每列约 $n$ 条支持，terminal 距离为 $O(QnK)$；
numeric LL 的协方差与求解为 $O(QnK^2+QK^3)$，离散输出另有类别域存储。
把所有缺失 cells 都查询时 $Q$ 可随 $NM$ 增长，整体推理不再由 backbone 的线性项决定。

\section{数学性质与实现必须守住的边界}
\label{sec:properties}

\begin{proposition}[Null 闭合]
在所有中间量有限时，每一层的 Null carrier 恒等于零。
\end{proposition}
\begin{proof}
初始 Null 为零。式~\eqref{eq:presence} 给出 $\rho(0)=0$，attention receiver gate 删除更新；
FFN gate 也为零，最后显式 Null write-back 再保持零。逐子层归纳即可。
投影 bias 无法绕过 projection 后的 gate。
\end{proof}

\begin{proposition}[固定证据下的 Query 插入不扰动]
固定 schema、$\V$、可见值、参数及 codebook，把额外 Null 地址改成 Query，已有 carriers 和预测不变。
\end{proposition}
\begin{proof}
新增 Query 只改变自身初值。Collect 与 row sources 只来自 $\V$；其他 receiver 的更新只依赖
自己的状态与这些 sources。逐层归纳得到原 Value、banks 与旧 Query 状态不变。
响应公式与支持集合也不变，因此 terminal 输出不变。
\end{proof}

这个结论不覆盖 Value$\to$Query（那会移除事实），也不覆盖新增可见测试 covariates。
训练期若启用随机 dropout 或遍历顺序相关的 codebook，则需要共同搬运随机 realization；本稿 dropout 为零。

\begin{proposition}[地址置换等变]
同时置换行列、schema、masks、values、codebook 和输出地址，不改变对应预测的值。
\end{proposition}
\begin{proof}
编码参数共享且无 learned row/column ID；固定映射下编码随地址移动。
列/行注意力的合法 sources 随同置换，softmax 加权求和不依赖排列顺序。
Inducing seeds 跨列共享，语义 bank 的 $k$ 轴保持固定。
逐层传播、响应、同列 LL/NW 均保持地址对应，得到结论。
\end{proof}

有限精度下按容差核对，而非默认 bitwise equality。任意重排 ordinal 类别而不携带声明次序，不属于同一输入语义。
只改变 $K$ 或 $S$ 的值也不是等变变换。

\subsection{压缩与信息限制}
Collect/read 一般不等价于 direct attention，即使 $S\ge N$ 也没有精确恢复保证。
固定两次单头 attention 的权重、暂忽略 residual/FFN 时，其行到行混合经过 $S$ 维槽轴，
混合矩阵的秩至多为 $S$；完整带非线性模型不能简单套用这个线性秩结论。
这个视角解释了压缩的取舍，却不能证明哪种结构预测更好。

同列 supports 的值在编码中可见，Query 值不可见；两类 Cell 的分布不同。
共享 $W$ 和直接距离要求模型学出可比较的表示，这需要训练而不是类型检查来保证。
同样，两个 Query seeds 相近、responses 尺度失控、softmax 过尖或接近均匀，都可能阻碍学习。
这些现象应测量，不靠静默加温度、detach 或隐藏 gate schedule 消除。

\section{实现者的模块边界与验收}
\label{sec:acceptance}

\subsection{建议的六个纯接口}
接口名称是实现建议，输入/输出和信息边界由本文定义：
\begin{longtable}{@{}p{37mm}p{113mm}@{}}\toprule
接口 & 合同 \\\midrule\endhead
\texttt{compile\_episode} & 只读 $E_{\rm in}$，返回 $H^{(0)}$、固定 masks、typed statistics、codebook。
不可读取 truth sidecar；无效输入给出明确状态。\\
\texttt{omab} & 输入 receiver、eligible sources、Null receiver mask；按第~\ref{sec:block} 节返回同形 receiver。
Source 非空判定与 head 共享。\\
\texttt{axis\_block} & 输入 $H^{(\ell)},\V$；返回完整 $H^{(\ell+1)}$。默认 collect/read/row，
direct 是独立模块配置。不得循环更新某列后让未更新列读取混合时间状态。\\
\texttt{response} & 从最终 $H$ 提取三类视图，返回 Query 和 support 的 $z$；同一响应公式，
不读取原始 truth。\\
\texttt{typed\_terminal} & 只读 target response、同列 support responses/values 与可见统计；
返回点预测/PMF/状态。没有持久参数，不写回 $H$。\\
\texttt{score} & 唯一接收 $Y^\star_\Q$ 的接口；按非空分支均值求和，返回 scalar loss 与类型统计。\\\bottomrule
\end{longtable}

Batch 可以是不同形状 episode 的列表。若 padding 成矩形，跨 episode attention 永不连通；
先逐 episode 平均 loss，再对 batch 平均。不能按不同表的 token 数重新加权目标。
Checkpoint 应包含实际权重、完整配置、code 版本、优化器/scheduler 状态与 RNG 状态；
具体仓库 ModelSpec 身份待实现时新建，不能借用旧结构 checkpoint 的身份。

\subsection{开始训练前应通过的检查}
这些是未来实现的验收标准，不是本稿已经完成的模型实验：
\begin{enumerate}
\item \textbf{参数合同。}全尺寸实现的参数张量名称/形状能映射到计数脚本，总量 54,071,520；
没有遗漏 bias，也没有未声明的 norm、projection 或 target embedding。
\item \textbf{数据隔离。}固定 forward input，任意替换 sidecar labels，predictions 不变；
更改 visible 值则允许变化。隐藏值不能影响均值、方差、类别域或 RNG 映射。
\item \textbf{零与空源。}有 bias、空列、空行、真实值零、很小非零 carrier、全 Null corner 均覆盖；
Null 输出严格零，空源 attention update 零，非 Null 的局部 FFN 不应被错误跳过。
\item \textbf{固定槽与分块。}$N\in\{2,8,127,128,129\}$ 均使用 $S=128$；
完整 sources 与不均匀 chunk 的归约结果一致。不能把 chunk 各自 softmax 后平均。
\item \textbf{置换与 Query 插入。}固定 codebook，置换地址或插入 Null$\to$Query 后核对旧预测；
同时检查梯度对应关系。新增 visible 事实不是这个测试的等价变换。
\item \textbf{Terminal 代数。}LL 的中心求解与带截距正规方程一致；公共 Feature 平移不改 LL/NW；
重复支持保留 multiplicity；PMF 含声明但未见类别且和为 1。
\item \textbf{梯度。}在小随机表上以 FP64 有限差分检查非退化的 $W,w_U$、Cell seeds、attention 与 LL 求解梯度。
检查 support 路径不是 detach；Feature 专属梯度应为零，不强求所有参数都非零梯度。
\item \textbf{训练链。}先在独立小表上确认有限 loss、optimizer 更新、重放和验证分离，再开始预算训练。
过拟合一张表只证明该训练链在该样例可拟合，不证明 ICL。
\end{enumerate}

建议 FP64 代数比较 $\texttt{atol}=10^{-9},\texttt{rtol}=10^{-7}$；
FP32 forward 比较 $10^{-5},10^{-4}$。这些是起始容差，必须同时报告最大误差与输入尺度，
不能只放宽阈值到测试通过。激活大到误差超界时，应先检查 conditioning 和运算顺序。

\subsection{本稿随附的轻量校验}
\texttt{design/parameter\_count.py} 检查张量与参数公式；
\texttt{design/verify\_algebra.py} 仅检查 LL、PMF、平移与流式归约的独立小算例。
二者都不是训练模型实现，也不代替上述端到端验收。
\texttt{design/defaults.json} 供实现者加载本稿数值；公式、初始化、sampler 与 prior 的语义仍在本文完整定义。

\section{未来论文的实验骨架与当前开放问题}
\label{sec:experiments}

本文可以直接生长为论文，但现阶段不填写虚构的结果表、SOTA 摘要或优势曲线。
未来结果章节应围绕表~\ref{tab:experiments} 的问题补充真实证据，方法章节在配置改变时同步修订。

\begin{table}[htbp]\centering\small
\caption{待完成实验，不含任何已有性能结果。}\label{tab:experiments}
\begin{tabularx}{\linewidth}{@{}lYY@{}}\toprule
问题 & 固定或对照 & 需要报告的量 \\\midrule
模型是否学会恢复 & 独立生成世界、不同 mask 与 seed & numeric 标准化误差；离散 NLL、准确率；失败比例 \\
是否迁移到监督 ICL & 固定预处理、数据划分与测试行语义 & 独立真实数据集结果；均值、区间与分任务分布 \\
Inducing 的取舍 & 同深度 direct；近似同参数 direct 分开 & 效果、参数、FLOPs 估计、显存、延迟 \\
小样本是否受损 & 固定模型、改变可见 $N$ & $N<S$ 与 $N\ge S$ 下的曲线；固定槽不变 \\
响应几何能否学习 & Cell-only、Unit 开关；Feature 等价控制 & entropy、$n_{\rm eff}$、响应尺度、梯度与条件数 \\
读出选择 & Gaussian 固定；单独比较 learned matching & LL/NW 同 estimator 对照；参数与训练预算 \\
规模能力 & 分别扫 $N,M,Q$、支持数与类别数 & 质量、显存、耗时与 OOM/数值失败边界 \\\bottomrule
\end{tabularx}
\end{table}

与 TabPFN-3 比较时应报告具体分类/回归 checkpoint、单 estimator 或 ensemble、预处理和总推理成本，
不能只比较名义参数量。TAR 初版 numeric 点预测不具备对方分布回归输出；分位数或 calibration 比较要等
对应输出接口存在后再开展。数据任务的合法可见信息与测试批语义应对齐，否则结论需限定到不同协议。

\subsection{哪些决定可以先固定，哪些结论不能提前写}
当前可实施的决定是完整 tensor program、54M 预算、单轴固定槽、Gaussian LL/NW 和 masked-cell loss。
接下来最可能推动修订的是响应尺度稳定性、Cell Query 与 Value 的表示差异、Feature 通道不可辨识、
synthetic prior 不足以及逐 Cell 状态的内存成本。
应以实验或推导定位问题，再修改对应接口；无需为了预见一切而在首版同时启用所有候选。

成功的第一轮实现应能准确回答：“运行的每个算子是否对应本文、数据是否正确隔离、有限差分与分块是否一致、
哪些数据形状可以运行、哪些统计表明模型在学习”。达到这些条件后，再谈训练预算扩大与能力比较。

\appendix
\section{完整 forward 伪代码}
\begin{enumerate}
\item 验证 schema、$\V,\Q$、可见数值及支持数；构造 truth-free $E_{\rm in}$。
\item 只用 $\V$ 拟合连续统计、创建 nominal codebook；按式~\eqref{eq:compiler} 构造 $H$、Null mask。
\item 对 $\ell=0,\ldots,L-1$：
先只从 $H$ 的 visible 地址为各列计算 $G_a$（空源覆盖为零）；
由同一个旧 $H$ 回读 $G_a$ 得到全部 $H^{\rm col}$；
再只读取 $H^{\rm col}$ 的 visible 地址完成 row OMAB，得到新 $H$。
Direct 配置以直接 column OMAB 替换 collect/read。
\item 从最终 $H$ 读取所有所需的 $c,U,F$，按式~\eqref{eq:response} 为 Query 与 supports 算 $z$。
\item 对每个 Query：提取同列完整支持；用式~\eqref{eq:matching} 计算权重；
按类型执行式~\eqref{eq:solve} 或式~\eqref{eq:pmf}，返回地址与状态。
\item 训练时 scorer 才读取 sidecar，按式~\eqref{eq:loss} 回传；推理直接返回预测。
\end{enumerate}

\section{文件职责与修订方法}
\texttt{MATHEMATICAL\_DESIGN.tex} 是共同数学维护稿；
\texttt{end-to-end-design.tex} 展开通用结构族及完整推导；
\texttt{model-design.tex} 是本份固定规模、默认值齐备的实现稿，也是未来论文骨架。
本稿为独立阅读重写必要定义，不复制完整历史附录；无须加载父文档的 aux 或 TeX。
只有共同底稿使用大写文件名，本文名称不编码模型代际或性能承诺。

若共同 response、OMAB、support、terminal 或 loss 改动，必须同步检查本稿相关公式、
图~\ref{fig:flow}--\ref{fig:cache}、defaults 和参数/代数校验。
现有 \texttt{check\_consistency.py} 检查共同底稿与通用端到端稿；它不自动证明本稿所有摘录一致。
本轮没有改变共同数学、旧结构合同、RecSys 设计或运行代码。

\begin{thebibliography}{9}
\bibitem{pfn3} Prior Labs Team. \emph{TabPFN-3: Technical Report}. arXiv:2605.13986v2, 2026-05-28.
\url{https://arxiv.org/abs/2605.13986v2}. 本稿参照其参数规模与附录 C 的架构配置。
\bibitem{pfncode} Prior Labs. \emph{TabPFN v3 architecture implementation}.
固定 commit \texttt{bfc3cbb9c5532e610ea01dbdee1f765057b5a8e3}.
\href{https://github.com/PriorLabs/TabPFN/blob/bfc3cbb9c5532e610ea01dbdee1f765057b5a8e3/src/tabpfn/architectures/tabpfn_v3.py}{官方代码快照}，2026-09-06 核对。
\end{thebibliography}
\end{document}
